\documentclass[10pt,a4paper]{amsart}
\usepackage[margin=19mm]{geometry}
\usepackage{amsmath,amssymb,mathtools}
\usepackage[scr=rsfs,cal=euler]{mathalfa}
\usepackage{xfrac}
\usepackage[unicode,hidelinks,bookmarks=false]{hyperref}
\newcommand{\ie}{i.e.\ }
\newcommand{\cf}{cf.\ }
\newcommand{\resp}{respectively\ }
\newcommand{\Subsection}{\subsection}
\providecommand{\nobreakdash}{\nobreak\hbox{-}}
\setlength{\emergencystretch}{2em}
\multlinegap0pt
\usepackage{bm}
\usepackage{bbm}
\usepackage{dsfont}
\usepackage{xspace}
\usepackage{cancel}
\usepackage{mathtools}
\usepackage{amsthm}
\makeatletter
\@ifundefined{theorem}{\newtheorem{theorem}{Theorem}}{}
\@ifundefined{proposition}{\newtheorem{proposition}[theorem]{Proposition}}{}
\makeatother
\def\gcd{\operatorname{gcd}}
\def\mod{\operatorname{mod}}
\title[Ranks of Elliptic Curves Twisted by Quadratic Forms]{Ranks of Elliptic Curves Twisted by Quadratic Forms}
\author{Mohammad H. Hamdar, Cihan Sabuncu}
\date{Source: arXiv:2607.13000v2 (CC BY 4.0)}
\begin{document}
\maketitle

\noindent\emph{Source and licence.} Ranks of Elliptic Curves Twisted by Quadratic Forms. Version arXiv:2607.13000v2 (CC BY 4.0), distributed under the Creative Commons Attribution 4.0 International licence, as recorded in the arXiv licence metadata for this exact version and shown on the abstract page https://arxiv.org/abs/2607.13000v2. Full rights notice: licenses/arxiv-2607.13000-CC-BY-4.0.txt.\par\smallskip

\noindent\textit{Editorial scope.} Counted unit: the Proposition whose source label is \texttt{G\_k(n) conditions} in the article (source0.tex lines 483--493, source label \texttt{G\_k(n) conditions}), delivered together with the article's own complete original proof of it (source0.tex lines 494--533, one \texttt{proof} environment of 40 lines). No mathematics is written, repaired or re-derived: statement and proof are the article's own text line for line. Required context retained uncounted: def\_of\_rho (lines 478--480). Every definition, display and label of those lines is retained unchanged. Omitted from this package and not redistributed: 1--477, 481--482, 534--822. Nothing inside a retained excerpt is omitted or altered: every retained body is delivered exactly as the article's lines are written, so no omission receipt of any kind accompanies this package. Citations: the retained excerpts carry no citation command at all, so no bibliography entry of the article is transcribed and no reference is redistributed. Reference adaptation: none was needed and none was made; every reference inside the delivered unit resolves inside it, so no unresolved reference or citation is produced. Printed slips kept literal: no printed word, symbol, hypothesis or number of the counted statement or of its proof was altered. Layout and class: the article's own class is \texttt{amsart} with packages including geometry, dsfont, amsmath, amsthm, amssymb, multirow, mathtools, times, inputenc, babel, indentfirst, graphicx, color, hyperref; the wrapper is the lane's pinned \texttt{amsart} editorial wrapper at 10pt on A4, which restates the theorem-like environments the retained text needs and adds the article's own macro definitions that the retained text needs (\texttt{\textbackslash gcd}, \texttt{\textbackslash mod}), with extra wrapper packages (bm, bbm, dsfont, xspace, cancel, mathtools). The delivered numbering is the unit's own. Assets: no retained span contains a figure environment, a graphics-inclusion command, a PDF-page inclusion, a table or a TikZ picture; no image file is copied into this package, the package declares no asset dependency and no image file is redistributed with it. Licence: version arXiv:2607.13000v2 of this article is distributed under the Creative Commons Attribution 4.0 International licence, as recorded in the arXiv licence metadata for this exact version and shown on the abstract page https://arxiv.org/abs/2607.13000v2. A full rights notice is delivered at licenses/arxiv-2607.13000-CC-BY-4.0.txt. Characters: every retained excerpt is pure ASCII, so the delivered engine, XeLaTeX with its default Latin Modern text font, reports no missing character.
\begin{equation}\label{def_of_rho}
\rho(n):=\rho_{(0,0)}(n)=\#\{ (u_1,u_2)\mod n : u_1^2+u_2^2\equiv 0 \mod n \}.    
\end{equation}

\begin{proposition}\label{G_k(n) conditions}
If $m$ and $n$ are relatively prime odd integers, then $$G_{(k_1,k_2)}(mn)=G_{(k_1,k_2)}(m)G_{(k_1,k_2)}(n).$$ Moreover, if $p^{\alpha}$ is the highest power of $p$  dividing $\gcd(k_1,k_2)$ (setting $\alpha=\infty$ if $(k_1,k_2)=(0,0)$), then we have for $\beta\geq 1$,
$$ G_{(k_1,k_2)}(p^{\beta}) = \begin{cases}
    0 & \text{if } \beta\leq \alpha \text{ is odd} \\
    p^{2(\beta-1)} (p^2 - \rho(p)) & \text{if } \beta\leq \alpha \text{ is even} \\
    p^{2\alpha+1} \left( \frac{(k_1^2+k_2^2)p^{-2\alpha}}{p}\right) & \text{if } \beta = \alpha + 1 \text{ is odd} \\ 
    -p^{2\alpha}\left( \frac{-1}{p} \right)(p\cdot \delta_p((k_1^2+k_2^2)p^{-2\alpha}) - 1)  & \text{if } \beta=\alpha + 1 \text{ is even} \\
    0 & \text{if } \beta\geq \alpha+2,
\end{cases} $$
where $\delta_p(n)$ is $1$ if $n\equiv 0 \mod p$ and $0$ otherwise, and $\rho(n)$ is as in \eqref{def_of_rho}.
\end{proposition}

\begin{proof}
    We first prove the multiplicativity. Since $\gcd(m,n)=1$, write $b_1 = m b_1'+ nb_1'' $ for $b_1' \mod n$ and $b_1'' \mod m$ by the Chinese remainder theorem. And similarly for $b_2=m b_2' + n b_2''$. We then have
    $$G_{(k_1,k_2)}(mn) = \sum_{\substack{b_1', b_2' \mod n \\ b_1'',b_2'' \mod m}}\left(\frac{(b_1' m+ b_1''n)^2+(b_2' m + b_2''n)^2}{mn} \right) e\left( \frac{k_1(b_1' m + b_1'' n) + k_2 (b_2' m + b_2'' n)}{mn} \right). $$
    Using the multiplicativity of the Kronecker symbol and the linearity in the exponential phase, we get that this is $G_{(k_1,k_2)}(m) G_{(k_1,k_2)}(n)$ as desired.

    We now prove the second part of the lemma. If $\beta=\alpha + 1$, then
    \begin{align*}
    \sum_{b_1,b_2 \mod p^{\beta}} \left( \frac{b_1^2+b_2^2}{p^{\beta}} \right) &e\left( \frac{k_1 b_1}{p^{\beta}} + \frac{k_2 b_2}{p^{\beta}}\right) \\
    &= \sum_{\ell_1 , \ell_2 \mod p} \left( \frac{\ell_1^2+\ell_2^2}{p^{\beta}} \right) \sum_{a_1 , a_2 \mod p^{\beta-1}} e\left( \frac{k_1 (a_1p + \ell_1) + k_2 (a_2p + \ell_2)}{p^{\beta}} \right) \\
    &= p^{2(\beta-1)} \sum_{\ell_1 , \ell_2 \mod p} \left( \frac{\ell_1^2+\ell_2^2}{p^{\beta}} \right) e\left( \frac{k_1 \ell_1 + k_2 \ell_2}{p^{\beta}} \right).
    \end{align*}
    If $\beta$ is even, then the sum above is
    $$ \sum_{\substack{\ell_1 , \ell_2 \mod p \\ \ell_1^2+\ell_2^2 \not\equiv 0 \mod p}} e\left( \frac{k_1 \ell_1 + k_2 \ell_2}{p^{\beta}} \right) = - \sum_{\substack{\ell_1 , \ell_2 \mod p \\ \ell_1^2+\ell_2^2 \equiv 0 \mod p}} e\left( \frac{k_1 \ell_1 + k_2 \ell_2}{p^{\beta}} \right) ,$$
    since the full sum is 0. We write $k_1'=k_1/p^{\alpha}$ and $k_2'= k_2/p^{\alpha}$. By orthogonality, we get
    $$ \sum_{\substack{\ell_1 , \ell_2 \mod p \\ \ell_1^2+\ell_2^2 \equiv 0 \mod p}} e\left( \frac{k_1 \ell_1 + k_2 \ell_2}{p^{\beta}} \right) = \frac{1}{p} \sum_{a\mod p}\sum_{\substack{\ell_1 , \ell_2 \mod p}} e\left( \frac{a(\ell_1^2+\ell_2^2) + k_1' \ell_1 + k_2' \ell_2}{p} \right) .$$
    This equals
    $$ \frac{1}{p}\sum_{a\mod p} \left( \sum_{\ell_1 \mod p} e\left( \frac{a\ell_1^2+ k_1'\ell_1}{p} \right) \right) \left(\sum_{\ell_2 \mod p} e\left( \frac{a\ell_2^2+ k_2'\ell_2}{p} \right) \right) . $$
    Note that the term $a=0$ gives $0$. By the theory of Gauss sums, for $a\not=0$, these sums over $\ell_1, \ell_2$ evaluate to
    $$ \left( \frac{a}{p} \right) \tau_p e\left( \frac{-\overline{4a}k_i'^2}{p} \right), $$
    where $\tau_p$ is the usual Gauss sum and so $\tau_p^2 = \left( \frac{-1}{p} \right)p$.
    Hence, we get the fourth case of the lemma, since we are left with
    $$ - p^{2(\beta-1)} \left( \frac{-1}{p} \right) \sum_{\substack{a \mod p \\ a\not= 0}}e\left( \frac{-\overline{4a}(k_1'^2+k_2'^2)}{p} \right) , $$
    and this sum is the Ramanujan sum.

    If $\beta$ is odd, we have
    $$ p^{2(\beta-1)} \sum_{\ell_1 , \ell_2 \mod p} \left( \frac{\ell_1^2+\ell_2^2}{p} \right) e\left( \frac{k_1 \ell_1 + k_2 \ell_2}{p^{\beta}} \right) . $$
    Similar to the even case, we can write the sum as
    $$ \sum_{\ell_1 , \ell_2 \mod p} \left( \frac{\ell_1^2+\ell_2^2}{p} \right) e\left( \frac{k_1' \ell_1 + k_2' \ell_2}{p} \right) .$$
    We rewrite the Legendre symbol as
    $$ \left( \frac{\ell_1^2+\ell_2^2}{p} \right) = \frac{1}{\tau_p} \sum_{\substack{a \mod p \\ a\not=0}} \left( \frac{a}{p}\right) e\left( \frac{a(\ell_1^2+\ell_2^2)}{p} \right) . $$
    Putting this in, similar to the even case, we get sums of the form
    $$ \sum_{\ell_i \mod p} e\left( \frac{a\ell_i^2 + k_i' \ell_i}{p} \right) = e\left( - \frac{\overline{4a}k_1^2}{p} \right) \sum_{\ell_1} e\left( \frac{a(\ell_i + \overline{2a} k_i')^2}{p} \right). $$
    Now changing the variable of summation $\ell_i + k_i' \overline{a} \to \ell_i$, the inner sum evaluates to $\left( \frac{a}{p} \right)\tau_p$ by the classical theory of Gauss sums. Putting everything back and using $\tau_p^2 = \left( \frac{-1}{p} \right)p $ again, yields
    $$ \sum_{\ell_1 , \ell_2 \mod p} \left( \frac{\ell_1^2+\ell_2^2}{p} \right) e\left( \frac{k_1' \ell_1 + k_2' \ell_2}{p} \right) = \frac{\left(\frac{-1}{p} \right) p}{\tau_p} \sum_{\substack{a \mod p \\ a\not=0}} \left( \frac{a}{p} \right) e\left( \frac{-\overline{4a}(k_1'^2+k_2'^2)}{p} \right) .$$
    Lastly, we do the change of variables $-\overline{4a} \to u$ which gives $\left( \frac{a}{p} \right) = \left( \frac{-\overline{4u}}{p} \right) = \left( \frac{-1}{p} \right)\left( \frac{u}{p} \right)$. Hence, we obtain
    $$ \frac{p}{\tau_p} \sum_{\substack{u \mod p \\ u\not=0}} \left(\frac{u}{p} \right) e\left( \frac{u(k_1'^2+k_2'^2)}{p} \right) .$$
    Evaluating this Gauss sum gives the third case of the lemma.

    The other cases are done in a similar manner.
\end{proof}

\end{document}
